% !TeX program = xelatex
% 第三篇论文：研究方法与实验协议（证据约束草稿）
% 本文件为独立可编译章节。并入学校学位论文模板时须删除导言区和document环境。
% 证据依据：prepare_dlpfc_10_sources.py / prepare_dlpfc_10_features.py /
% mvp_fixed15_dlpfc_10_sections.py / mvp_fixed15_block_swap_151508.py /
% cellcharter_151507_baseline_v032.py，以及已完成的Windows终端审计。
% 文献正式引文、旧开发两片的来源差异仍待最终核实，勿自行虚构。
\documentclass[UTF8,a4paper,12pt]{ctexart}
\usepackage[margin=2.5cm]{geometry}
\usepackage{amsmath,amssymb,booktabs,tabularx}
\usepackage{tikz}
\usetikzlibrary{arrows.meta,positioning}
\usepackage{hyperref}
\hypersetup{hidelinks}
\setlength{\parindent}{2em}
\setlength{\parskip}{0.3em}
\renewcommand{\arraystretch}{1.18}
\setcounter{section}{2}

\begin{document}
\section{研究方法与实验协议}
\label{sec:method}

\subsection{研究目标与实验范围}
本研究以已公开的CellCharter多尺度空间表征与无监督聚类流程为研究基础，考察不同邻域尺度特征块组合对人脑背外侧前额叶皮层（dorsolateral prefrontal cortex, DLPFC）空间域划分及组织层间边界检出的影响。研究不以提出新的神经网络为前提，而是在既有表征上构建\textbf{输入维度相同、聚类配置相同}的A/B对照，重点分析整体聚类指标与局部边界表现之间是否存在权衡。

研究涉及12张DLPFC切片，按组织来源划分为3位供体，每位供体4张切片。151507、151508为已有开发性观察，151509--151510和151669--151676为随后扩展的10张切片。所有切片均为技术意义上的样本切片，而不是12个独立供体。后续统计以切片结果进行配对描述，并明确仅有3位生物学独立供体。

\subsection{原始数据、坐标及参考标签}
\label{subsec:source}
每张切片的基因表达矩阵通过10x格式的 \texttt{filtered\_feature\_bc\_matrix.h5} 读取。与之配套的 \texttt{spots\_aligned.csv} 按表达矩阵中的barcode逐行对齐，包含Visium阵列坐标 \texttt{array\_row}、\texttt{array\_col} 及组织区域标记。表达矩阵行序与空间位置表的barcode要求完全一致，且不能存在重复barcode或重复阵列坐标；用于聚类的spot均须满足 \texttt{in\_tissue=1}。

后续扩展的10张切片由归档脚本 \texttt{prepare\_dlpfc\_10\_sources.py} 对照原始H5、位置表和SEDR来源的 \texttt{metadata\_source.tsv} 生成标准化表格。其中用于评价的参考层标签来自元数据列 \texttt{layer\_guess}，并与 \texttt{layer\_guess\_reordered} 进行一致性检查。标签名如\texttt{Layer1}被统一为\texttt{Layer 1}，未标注spot保留在表达及空间图中，但不进入有监督评价。\textbf{这些标签是既有元数据参考标注，并非本研究重新人工判读的组织学真值。}

旧开发切片151507、151508使用此前已准备的参考标注文件，不能在未核对原始来源与判读过程前断言其与后续10片完全同源。该参考标注来源差异在跨切片解释时保留为潜在混杂因素。

\subsection{表达预处理与scVI低维表示}
\label{subsec:scvi}
预处理依据归档运行脚本：首先使基因名唯一；分别以 \texttt{min\_counts=3} 过滤基因和spot，检验spot数量与barcode顺序未发生变化；随后在 \texttt{counts} 层保留原始计数副本。另对主表达矩阵使用总量归一化（\texttt{target\_sum=1e6}）和对数变换，并基于 \texttt{counts} 层使用 \texttt{seurat\_v3} 选出5000个高变基因（highly variable genes, HVGs）。

通过 \texttt{scvi.model.SCVI.setup\_anndata(..., layer="counts")} 设置模型输入，构建潜变量维数为5的scVI模型。源代码中的训练设置为最大400 epoch、启用提前停止（early stopping）、AdamW优化器及学习率平台期衰减，随机种子设为22911。后续10张切片的脚本要求使用CUDA训练，不允许静默退回CPU。模型训练得到每张切片的低维表示
\begin{equation}
  Z = \left[z_1^\top,\ldots,z_N^\top\right]^\top
      \in\mathbb{R}^{N\times 5},
\end{equation}
其中$N$为对应切片纳入表达矩阵的spot数。scVI的输入为 \texttt{counts} 层；归一化和对数变换及HVG处理按上述实际代码执行，不能将归一化后的矩阵误称为直接用于scVI拟合的计数输入。

各切片独立进行特征准备；已有Windows日志和缓存审计显示新增10张切片的5维表示与后续25维缓存已经生成并验证。训练最大epoch是上限，\textbf{不意味着各切片都训练到400 epoch或均已严格收敛}。

\subsection{空间邻接图与CellCharter多尺度聚合}
\label{subsec:aggregate}
根据Visium阵列坐标建立无向空间邻接图$G=(V,E)$。脚本通过 \texttt{squidpy.gr.spatial\_neighbors} 在六邻接网格上建图，并验证边满足阵列坐标偏移
\begin{equation}
  (|\Delta r|,|\Delta c|)\in\{(0,2),(1,1)\}.
  \label{eq:hex}
\end{equation}
同时检查邻接矩阵对称、每个节点邻居数不超过6，以及实际无向边数与预先核定的样本值一致。这里的图由物理坐标建立，\textbf{并非根据基因表达距离构建的k近邻图}。

将scVI五维嵌入存入 \texttt{X\_scVI}，调用
\texttt{cc.gr.aggregate\_neighbors(adata, n\_layers=4, use\_rep="X\_scVI")}，
生成CellCharter多尺度聚合表示 \texttt{X\_cellcharter}。本研究沿用此软件实现，不人为宣称其四个聚合层均采用简单均值或具有尚未核实的精确权重公式。程序保存的矩阵维数为
\begin{equation}
  H=[L_0\,|\,L_1\,|\,L_2\,|\,L_3\,|\,L_4]
    \in\mathbb{R}^{N\times25},
  \quad L_{\ell}\in\mathbb{R}^{N\times5}.
  \label{eq:multi}
\end{equation}
式中$L_0$表示未聚合的局部表示块；$L_1$至$L_4$对应CellCharter API返回的逐层特征块位置，\textbf{不将索引名称自行推广为已验证的生物学影响半径}。

生成的特征以 \texttt{scvi\_latent.npy} 与 \texttt{cellcharter\_aggregated.npy} 保存，并记录 \texttt{spot\_order.csv}、配置、状态及SHA256校验值。冻结比较只读取已经生成的25维缓存，不重新训练scVI。

\subsection{等维特征组合与配对GMM聚类}
\label{subsec:ab}
为隔离“选择哪些邻域特征块”这一对照因素，同时避免比较15维与25维输入造成维数混杂，固定两组15维特征：
\begin{align}
  X^{(A)} &= [L_0\,|\,L_1\,|\,L_2]\in\mathbb{R}^{N\times15},\\
  X^{(B)} &= [L_0\,|\,L_1\,|\,L_4]\in\mathbb{R}^{N\times15}.
\end{align}
两组仅替换第三个5维块，均通过从相同25维缓存中按列截取得到。这里的B是\textbf{待检验的固定候选配置}，不是已证实优于A的新算法。两组分别拟合GMM，而非把A的预测结果作为B的初始化；每次拟合前均重置NumPy及PyTorch的全局种子为22911。

基于 \texttt{CellCharter~0.3.2} 的 \texttt{cc.tl.Cluster} 接口，两个条件采用相同设置：全协方差（full covariance）、K-means初始化、收敛容差$10^{-3}$、协方差正则$10^{-6}$、CPU训练、最大100 epoch、单设备。聚类数$K$取参考层类别数：供体D2四张切片$K=5$，其余8张切片$K=7$，属于\textit{Oracle-K}协议。这一做法使用了参考标注的类别数量，不能解释为完全无标签选择K；但在扩展10片的冻结A/B脚本中，\textbf{具体spot层标签是在两组GMM预测保存后才读取用于评价}，未参与GMM参数拟合。

在实际版本依赖下，\texttt{random\_state}向GMM内部的传递细节尚不能仅凭外层接口保证所有随机路径严格一致。本研究已记录版本及手动设置全局种子的措施，\textbf{不声称完成跨随机种子的稳定性检验}。

\subsection{实验技术路线}
图\ref{fig:pipeline}根据已审计源代码绘制，为研究流程示意而非训练结果图。参考标签仅进入聚类数Oracle-K设定及预测后的评价环节。

\begin{figure}[htbp]
\centering
\begin{tikzpicture}[
  font=\small,
  node distance=5.5mm,
  stage/.style={draw,rounded corners=2pt,align=center,text width=10.6cm,
                minimum height=0.72cm,inner sep=5pt},
  arrow/.style={-{Latex[length=2mm]},thick},
]
\node[stage] (raw) {10x表达H5 + Visium阵列坐标 + 元数据来源核对\\barcode/spot顺序、来源哈希检查};
\node[stage,below=of raw] (pre) {预处理：过滤计数阈值3；保存counts；5000 HVG（seurat\_v3）};
\node[stage,below=of pre] (scvi) {scVI：5维潜变量；种子22911；CUDA训练；最多400 epoch};
\node[stage,below=of scvi] (graph) {六邻接空间图（由阵列坐标构建）+ CellCharter聚合4层};
\node[stage,below=of graph] (cache) {冻结并审计25维缓存：$[L_0|L_1|L_2|L_3|L_4]$};
\node[stage,below=of cache] (ab) {公平A/B特征对照：A=$[L_0|L_1|L_2]$，B=$[L_0|L_1|L_4]$\\两者均15维，独立同配置GMM，Oracle-K};
\node[stage,below=of ab] (eval) {预测后读取参考层标签：ARI、NMI、边界F1与TP/FP/FN\\12片配对描述、供体分组与典型层间边界分析};
\foreach \a/\b in {raw/pre,pre/scvi,scvi/graph,graph/cache,cache/ab,ab/eval}
  \draw[arrow] (\a.south) -- (\b.north);
\end{tikzpicture}
\caption{基于实际脚本的CellCharter多尺度特征对照流程图。原始参考标签只用于Oracle-K设置及聚类后评价；该图为方法流程示意，不代表原创模型结构。}
\label{fig:pipeline}
\end{figure}

\subsection{评价指标与分层边界分析}
\label{subsec:evaluation}
对每个有参考标注的spot，将其预测聚类编号与参考标签进行比较，计算调整兰德指数（ARI）与归一化互信息（NMI）。这些指标对预测聚类编号置换不敏感，不能根据相同数值簇ID推断两个配置对应相同生物学层。

将无向六邻接边仅计数一次，且只保留两端均存在参考层标注的边。令$y_i$与$\hat y_i$分别为参考层标签和预测簇编号，则对于合格边$(i,j)$，
\begin{align}
 b_{ij}^{\mathrm{ref}}&=\mathbb{I}(y_i\neq y_j),\\
 b_{ij}^{\mathrm{pred}}&=\mathbb{I}(\hat y_i\neq\hat y_j).
\end{align}
比较两个二值边界标识后得到TP、FP、FN：
\begin{align}
 P_{\mathrm{boundary}} &= \frac{TP}{TP+FP},\\
 R_{\mathrm{boundary}} &= \frac{TP}{TP+FN},\\
 F_{1,\mathrm{boundary}}&=\frac{2TP}{2TP+FP+FN}.
\end{align}
遇到零分母时按分析脚本返回0。对真实标签不同的参考边，根据两端标签构成的无序组织层对分组，分别计算边界检出率$TP_g/N_g$。这一分层值本质上是\textbf{真实边界的召回率}，不能替代每个层对的precision或归因于该层对的FP。

两张说明性切片151671与151676是在看到整体边界F1方向后选取，应用相同的空间邻接边定义与层对统计，仅用于解释已存在的正负案例，不视作事先随机选择的独立验证样本。全12片汇总值是切片的描述性比较，3位供体才是独立生物学来源，不进行以spot数或边数充当$n$的显著性分析。

\subsection{可复现性与研究边界}
本研究已有来源转换记录、5维/25维特征缓存、A/B预测CSV、结果JSON、元数据校验和12片汇总CSV。新增10片的barcode对齐、ARI/NMI以及边界TP/FP/FN和F1已在Windows执行独立重算审计并通过；旧两片开发切片的等价范围独立复核仍需谨慎界定。相关文件的实际本地存在性、哈希以及图表数值必须以最终原始档案为准，而不能仅根据文字记录推断。

本章研究方法是对成熟CellCharter/scVI流程的复现和有控制的特征块比较；尚没有证据支持将固定$L_4$替换$L_2$命名为一个新的算法或已验证的原创结构。本方法在真实分析中依赖Oracle-K、有限独立供体以及非完全一致来源的参考标签，论文结果应同时报告这些前提与局限。

\section*{定稿核查备忘（不属于正式论文正文）}
\begin{itemize}
\item 待补真实参考文献：DLPFC数据集、CellCharter、scVI、SEDR标签来源；尚未核实DOI，不在这里编造。
\item 已检查主要参数的来源脚本：\texttt{prepare\_dlpfc\_10\_sources.py}、\texttt{prepare\_dlpfc\_10\_features.py}、\texttt{mvp\_fixed15\_dlpfc\_10\_sections.py}、\texttt{mvp\_fixed15\_block\_swap\_151508.py}。
\item 151507/151508和新增10片的标签来源与运行协议差异需要最终逐项核对；本稿没有写成完全统一盲测。
\item 独立编译使用XeLaTeX；并入学校模板时删除本文件导言区和document环境，调整编号及图注。
\end{itemize}
\end{document}
